\documentclass{situationist-preprint}

\title{Why Explained Jokes Are Not Funny}
\msid{AES-13-b}
\msfamily{AES-13}
\msdiscipline{Aesthetics}
\mstype{Consequence}
\mskeywords{humor; incongruity; explanation; amusement}

\begin{document}
\maketitle
\vspace{-1em}

\begin{abstract}
On the account of AES-13-a, a joke is funny when it resolves an incongruity only partially. We derive from this why an explained joke ceases to amuse: explanation resolves the residual incongruity itself, by naming it as a device, and so removes the condition of amusement.
\end{abstract}

\section{The derivation}

By Definition~2.2 of AES-13-a, an incongruity is partially resolved when the reading that makes the incongruous element fit is itself in tension with what the audience knows. By Principle~P1, amusement is the pleasure of grasping the resolution while continuing to perceive what remains incongruous in it. Theorem~2.3 states that a joke whose incongruity is resolved entirely is not funny.

An explanation of a joke does two things. It makes the resolving reading explicit, and it accounts for the residual tension by classifying it: \enquote{the second sentence exploits an ambiguity of attachment in the first}. The classification gives the remaining incongruity a place in a coherent account of what the joke is doing. What was perceived as incongruous is now understood as the intended mechanism, and understanding it as such is a resolution of it. The joke as explained is therefore an incongruity resolved entirely, and by Theorem~2.3 it is not funny.

\begin{proposition}\label{p:expl}
An explanation of a joke that is grasped by its audience removes the joke's capacity to amuse that audience on that occasion.
\end{proposition}

The account predicts two refinements. First, the effect is distinct from familiarity: a joke heard often loses surprise, but an unexplained familiar joke can still amuse, because its residual incongruity is still perceived; an explained joke cannot, however fresh. Second, the effect concerns the token joke, not humor in general. A theory of humor explains types of joke, and grasping it does not resolve the incongruity of an unexplained joke one has not yet analyzed, which is why theorists of humor continue to laugh \citep{Carroll2014,HurleyDennettAdams2011}.

Explanation is fatal to a joke because it completes what the joke must leave incomplete.

\begin{thebibliography}{9}\small
\bibitem[Carroll(2014)]{Carroll2014} Carroll, N. (2014). \emph{Humour: A Very Short Introduction}. Oxford University Press.
\bibitem[Hurley et~al.(2011)]{HurleyDennettAdams2011} Hurley, M.~M., Dennett, D.~C. and Adams, R.~B. (2011). \emph{Inside Jokes: Using Humor to Reverse-Engineer the Mind}. MIT Press.
\end{thebibliography}

\end{document}
